%
% Chapitre "Structure d'un rapport technique"
%

\chapter{Besoins et objectifs}
\label{s:objectifs}
\section{Analyse des besoins}
Avant de se lancer dans la conception du projet et dans le but d'offrir au client le meilleur produit possible,
plusieurs besoin principaux ont été cernés. Premièrement, le système conceptualisé devra fonctionner de manière
complétement autonome. En effet, notre mandat exige que notre système devra être en mesure d'effectuer la prise
de données de manière complètement autonome pour une durée d'un an. À noter qu'il sera accessible selement pendant
2 semaine lors du mois de juin.

\bigskip
Deuxièmement, le système doit être robustre et doit être capable de supporter de gros écarts de température.
En effet, puisque ce système sera déployé dans l'artique québecois sur une période d'un an, il pourrait s'exposer
à des écarts de température pouvant varier entre -20°C et 20°C. De plus, le système devra être robustre et solide
afin d'être en mesure de fonctionner même si'il peut être recouvert par plus de 2 mètres de neige pendant une période
de 8 mois.

\bigskip
Troisièmement, le système devra limiter ses interactions avec la faune. Puisque ce dispositif sera potentiellement
installé dans des zones protégées, il devra assurer une mesure passive des données. En d'autres mots, il ne devra pas
piéger la faune afin de receuillir ses données. De plus\dots

\section{Analyse des objectif}
\section{Hiérarchisation des objects}
